% =============================================================================
% example-hr-job-description.tex
% Style: latex-docs-hr
% Build: pdflatex via latexmk; minted-aware with listings fallback.
% =============================================================================
\documentclass[11pt]{article}
\usepackage{latex-docs-hr}

\setDocTitle{Senior Documentation Engineer}
\setDocSubtitle{Engineering -- Documentation Tooling}
\setDocVersion{v1.0}
\setDocStatus{Open}
\setDocOwner{People Operations}
\setDocAuthor{Jordan Suber}
\setDocDate{May 8, 2026}
\setDocSummary{%
  Reference job description demonstrating the HR family module: JD
  header block, structured section macros, an applicant-record
  excerpt, and a performance rubric template. Personally identifiable
  information is rendered through the PII macro.%
}

\begin{document}
\makedoctitle
\makedocsummary

\begin{jobdescription}{Senior Documentation Engineer}{Engineering}%
                     {Director, Engineering}{Lithia Springs, GA / Remote (US)}
\end{jobdescription}

\subsection*{About the role}
We are looking for a senior engineer who treats documentation as a
first-class product. You will own the modular LaTeX style hierarchy,
authoring and reviewing the canonical, zero-warning examples that every
team builds on.

\Responsibilities
\begin{itemize}
  \item Maintain the \texttt{latex-docs} module hierarchy: parent,
        family modules, and technical subtypes.
  \item Author canonical, zero-warning example documents for each
        concrete module.
  \item Review documentation pull requests for adherence to the house
        style and the modular architecture.
  \item Partner with the Compliance and Legal teams to keep the
        redaction, retention, and signature workflows aligned with
        policy.
\end{itemize}

\Qualifications
\begin{itemize}
  \item 5+ years of professional experience producing technical
        documentation at scale.
  \item Strong LaTeX and \texttt{tcolorbox} fluency; comfortable with
        \texttt{titlesec}, \texttt{cleveref}, and \texttt{longtable}
        idioms.
  \item Demonstrable taste for typography and layout; an instinct for
        the difference between behaviour that belongs in the parent and
        behaviour that belongs in a child module.
\end{itemize}

\NiceToHaves
\begin{itemize}
  \item Experience with PlantUML and a CI/CD pipeline for documents.
  \item Exposure to enterprise architecture frameworks (Views and
        Beyond, ISO/IEC/IEEE 42010, TOGAF).
\end{itemize}

\Compensation
\begin{itemize}
  \item Competitive base; equity participation; comprehensive medical,
        dental, and vision.
  \item Annual learning and conference budget.
\end{itemize}

\HowToApply
\begin{itemize}
  \item Submit a short cover note and a representative documentation
        sample to \PII{careers@example.com}.
  \item Mention the job code \textbf{DOCENG-26-007}.
\end{itemize}

\section*{Hiring pipeline (sample)}
\begin{applicantrecord}{Candidate \#A2026-014, \PII{name redacted}}%
                       {DOCENG-26-007}
  \Stage{2026-04-22}{Application received}{Resume attached; portfolio link
         \PII{redacted}.}
  \Stage{2026-04-25}{Phone screen}{30-min screen with hiring manager;
         strong fit on tooling and style sense.}
  \Stage{2026-05-02}{Take-home}{Style-module migration exercise; output
         compiled with zero warnings on first attempt.}
  \Stage{2026-05-08}{On-site loop}{Four 45-min interviews; debrief
         scheduled for 2026-05-09.}
\end{applicantrecord}

\section*{Performance review template (use after 90 days)}

\begin{performancerubric}
  \Rubric{Technical fluency in LaTeX and tcolorbox idioms}{4}{Authored two
          new technical subtypes; both compile with zero warnings.}
  \Rubric{Architectural judgement (parent vs.\ child placement)}{4}%
         {Correctly resisted hoisting subtype-specific behaviour into
          the parent during the Q3 review.}
  \Rubric{Clarity of authored documentation}{3}{Examples are concise;
          some additional cross-references would help reviewers.}
  \Rubric{Cross-team partnership (Compliance, Legal, Engineering)}{4}%
         {Effective working relationships established within the first
          60 days.}
\end{performancerubric}

\end{document}



